\subsection{The geometric domains}

For a smooth projective variety $V/K$, write $\mathsf V(V)$ for
\[
 \CH^r(V)_{\Q,\geomhom}\subseteq\ker\ob^r_{V,K}
 \quad\text{for every }r.
\]
Here and below, the local assertions use the companion results stated in
\secref{sec:analytic}. A number-field variety has a \emph{full height
domain} if both complete groups in \eqref{eq:height-domain} are in the
domain of the companion normalized pairing for every $r$.

\begin{theorem}[Geometric domains]\label{thm:geometric-domains}
Let $K/\Q_p$ be finite. Each of the following smooth projective
geometrically integral varieties has property $\mathsf V$.
\begin{enumerate}[label=\textup{(\roman*)}]
\item A variety of dimension at most three; a geometrically uniruled
fourfold; or a variety of dimension at most five whose universal
geometric rational $\CH_0$ is supported in dimension at most three.
\item A product
\[
 T\times\prod_a C_a\times\prod_b S_b,
\]
where the $C_a$ are arbitrary smooth projective curves, the $S_b$ are
arbitrary smooth projective surfaces, and $T$ has potentially good
projective reduction. The coefficient $T$ may be a point.
\item A variety whose rational Chow motive, after a finite extension,
is a retract of a finite direct sum of Tate twists of the varieties in
\textup{(ii)}. In particular this includes products of point Hilbert
schemes of arbitrary surfaces with curves and the allowed coefficient
$T$, as well as the surface-moduli motives specified in
\secref{sec:products}.
\item A variety in one of the full-degree Chow-niveau, complete-intersection,
or flat-fibration cases of \secref{sec:niveau}, or satisfying an
all-degree numerical model criterion of \secref{sec:models}.
\end{enumerate}
For every specified geometrically homologically trivial cycle, property
$\mathsf V$ supplies a form satisfying \eqref{eq:strong-G} on its exact
support complement. For a number-field variety to which one of these
criteria applies at every finite completion, the normalized pairing has
a full height domain.
\end{theorem}

The proofs of (i), (ii)--(iii), and (iv) occupy
\secref{sec:endpoints}, \secref{sec:products}, and
\secref{sec:niveau}--\secref{sec:models}, respectively. The global
assertion follows from \lemref{lem:two-degrees}. Geometric constructions
can be base changed to all completions; a reduction condition on $T$ is
imposed separately at each finite place.

The distinction between a local degree range and a global height range
is useful even when all degrees are not covered. If $S\subseteq
\{1,\ldots,n\}$ is the set of local degrees established at every finite
place, then the full height pairs occur in
\begin{equation}\label{eq:range-intersection}
 S\cap(n+1-S).
\end{equation}
For example, small diagonal support gives local degrees $r\leq2$ or
$r\geq n-2$; for $n\geq6$ the corresponding full height degrees are
$r\leq2$ or $r\geq n-1$.

On an arbitrary smooth projective variety, the same formalism applies
to subgroups. In particular algebraically trivial cycles lie in the
local obstruction kernel in every degree. This gives the pairing on
the two algebraically trivial groups described in
\secref{sec:cycle-domains}.
